\section{Implementation}
\label{sec:implementation}

SAKI is implemented as an external Python controller plus two experiment
harnesses around unmodified \rocksdb instances. The implementation separates
the paper's two control questions and realizes the conserved rank-and-assign
framework from Section~\ref{sec:saki-design}. The epoch harness asks whether a
fixed budget can be placed better between workload phases; the continuous
harness asks whether the same placement loop can actuate a live RocksDB rate
limiter while tenants keep running. In both paths, the controller chooses high,
mid, and low budget assignments, the harness applies them through ordinary
RocksDB-facing APIs, and the measurement path records per-tenant behavior at
the same granularity used by the controller.

\subsection{Epoch Harness}

The epoch-level path uses \texttt{db\_bench} as the tenant process. A trial
creates one database directory per tenant, runs a prefill phase, and then
executes a sequence of workload epochs. Each epoch assigns tenants to high,
mid, and low demand tiers, launches one \texttt{db\_bench} process per tenant,
parses the logs, and uses the resulting observations for the next epoch's
allocation decision. Because every epoch is a fresh invocation, this harness
can vary both \texttt{max\_background\_jobs} and
\texttt{rate\_limiter\_bytes\_per\_sec} at epoch boundaries. The driver checks
that the resulting job and rate-limit settings match the static fair-share
budget before launch.

The parser extracts per-tenant runtime, exit status, operation rate, latency
histograms where available, compaction read/write bytes, flush bytes, keys and
bytes written, and compaction CPU time. It stores these rows with the epoch,
workload tier, assigned budget tier, background-job count, rate limit, and
workload parameters. This path gives clean mechanism evidence for budget
placement under drift because all policies run comparable foreground work under
the same aggregate budget. It does not by itself prove in-process runtime
control; that is the role of the continuous harness.

\subsection{Continuous Harness}

The continuous path is a custom C++ embedded RocksDB harness. Each tenant is a
long-running process linked against RocksDB. The Python driver writes one
control file per tenant, advances through fixed-length windows, computes the
current workload tier and policy assignment, validates the aggregate byte
budget, and atomically updates the control file. Tenant processes keep running
while their offered write/read rates, hot-key range, reported workload tier,
assigned budget tier, and byte budget are refreshed. Before each refresh the
driver checks that the next assignment preserves the aggregate byte budget.
Control files are written atomically, and a tenant keeps using its last
installed budget until it observes a newer value. Thus a missed or delayed
update does not create extra host budget in the harness; a production
controller would additionally fall back to all-mid allocation when tenant
inventory or metrics are incomplete.

Runtime actuation uses RocksDB's rate limiter. At startup, each tenant creates
a \texttt{NewGenericRateLimiter} with the initial budget and installs it before
opening the database. During the run, a control thread rereads the budget and
calls \ratecall on the same rate-limiter object. This is the main continuous
actuation path. Background-job count is a launch-time option and remains fixed
for the tenant process. Diagnostic userspace throttling exists in the harness
but is not the mechanism used for the reported continuous result.

The rate limiter is a RocksDB-facing background-I/O control rather than an
application-level admission throttle. In the reported setup it is installed as
the database rate limiter and therefore shapes the I/O issued by flush and
compaction work that uses RocksDB's rate-limited paths. SAKI does not directly
sleep foreground write calls or read calls. Foreground write latency changes
indirectly because background service changes the accumulation and draining of
LSM maintenance work. This distinction makes the continuous result a
maintenance-placement result rather than a foreground-throttling result:
SAKI changes the maintenance budget available to each tenant, not the offered
foreground workload. In the five main continuous trials, rate-limiter,
compaction, and flush byte deltas are larger in high-budget windows than in
low-budget windows. As an internal-validity check on this actuation path, we
also ran a homogeneous four-tenant static microstudy that sweeps the per-tenant
limiter budget \(B \in \{1,3,6,11\}\)\,MB/s with three seeds per point.
Dropping the first window, measured rate-limiter traffic is monotone in \(B\):
1.00, 2.96, 5.70, and 7.28\,MB/s. Compact-plus-flush output moves with the
limiter (coverage ratio 0.99--1.02); the 11\,MB/s point saturates at about
7.3\,MB/s because this workload cannot generate more maintenance I/O. A
two-seed no-rate-limiter sanity run reports zero limiter bytes while
background maintenance continues near 8\,MB/s. The microstudy shows that the
continuous knob controls background-maintenance I/O rather than foreground
throttling.

The embedded harness records per-window write and read operation counts,
throughput, and latency percentiles. A custom \texttt{EventListener} accumulates
compaction input/output bytes, compaction CPU time, compaction count, flush
bytes, and flush count. At window boundaries the harness emits deltas, together
with RocksDB properties such as pending compaction bytes, L0 files, memtable
bytes, delayed-write rate, and write-stop status. These measurements feed both
the online score and the post-run analysis.

\subsection{Analysis}

The analysis scripts consume JSON payloads produced by the drivers. For
continuous runs, the analyzer merges tenant metrics, allocation history,
process exit codes, commands, run parameters, and summary counters. It computes
write P99/P999 and throughput for high-pressure tenants, total throughput,
compaction output bytes, bytes/write, high-set overlap, adaptation lag, and
failed-tenant counts. The same pipeline emits the paper-table summaries and the
time-series data used for the evaluation figures. This keeps the manuscript
numbers tied to generated artifacts rather than manual transcription.
